\clearpage
\SetRegion{3.1 行政区划}{violet}{}
\phantomsection\label{ch3-sec31}
\pdfbookmark[1]{3.1 全球行政区划影响力排行榜}{sec31}
\section*{3.1 全球行政区划影响力排行榜}
\noindent\includegraphics[width=\linewidth]{figures/new/nav_31.pdf}\par\vspace{3pt}
加利福尼亚州、英格兰与纽约州列前三，榜首分值约为次席的1.40倍。前十名来自7个国家，完整Top 30涉及11个国家。

\plate{figures/31/01_global.pdf}{图3-1　全球行政区影响力 Top 30：名次、国别与分值}{来源：固定地区榜 \href{https://github.com/X-lab2017/2026-Global-Open-Source-Development-Report/blob/c51e3441361407b68f10b339d97eff288319b503/data/processed/chapter3/table-3-1-1_global_division_openrank_top30.csv}{[3.1-S1]}；统计窗口待核。州、省、构成国等单元不完全等价；条长以榜首为100\%，不代表全球份额。}

地区榜把宏观分布落实到具体节点。单个地区位居前列，与一国拥有多个入榜地区，是两种不同结构；后续按席位、位置和省市组合分别观察。

\clearpage
\pdfbookmark[2]{3.1.1 国别席位：覆盖与集中并存}{figure2}
\subsection*{3.1.1 国别席位：覆盖与集中并存}
Top 100覆盖32个国家。美国25席、印度12席、巴西9席、德国8席、中国7席，前五国合计61席；其余27国占39席。

\plate{figures/31/02_seats.pdf}{图3-2　全球 Top 100 行政区的国别席位图谱}{来源：固定席位汇总 \href{https://github.com/X-lab2017/2026-Global-Open-Source-Development-Report/blob/c51e3441361407b68f10b339d97eff288319b503/data/processed/chapter3/table-3-1-2_division_country_entry_counts.csv}{[3.1-S2]}；统计窗口待核。每格1席，代码为国家，按国家连续排列；不对应具体榜次或影响分值。}

19个国家各有1席。等面积方格呈现榜单覆盖，不能替代分值比较。各国行政区的数量和粒度不同，席位多也不直接等于国家影响力更强。

\clearpage
\pdfbookmark[2]{国别名次：席位相同，位置不同}{figure3}
\subsection*{国别名次：席位相同，位置不同}
美国在Top 30中有11个地区，其中3个进入前四；英国仅英格兰进入这一范围，但居第2。入榜数量与高位节点应联读。

\plate{figures/31/03_footprints.pdf}{图3-3　国别名次分布：由地区榜读取多节点结构}{来源：固定地区榜及席位汇总 \href{https://github.com/X-lab2017/2026-Global-Open-Source-Development-Report/blob/c51e3441361407b68f10b339d97eff288319b503/data/processed/chapter3/table-3-1-1_global_division_openrank_top30.csv}{[3.1-S1、S2]}；统计窗口待核。格点对应Top 30名次，右侧为两种范围内的席位数，非分值之和。}

中国与德国各有4席：中国的北京、上海、广东、台湾列第14、16、25、26，德国的四席列第9、15、22、28。印度在Top 100汇总中有12席，Top 30明细中有2席；缺少第31—100名明细，不能还原其余节点的位置。

\clearpage
\pdfbookmark[2]{3.1.2 参与与影响：同样本的两种位次}{figure4}
\subsection*{3.1.2 参与与影响：同样本的两种位次}
同一组30个行政区中，华盛顿州的规模位次为13、影响位次为4；卡纳塔克邦分别为2、8。参与规模与网络位置并不重合。

\plate{figures/31/04_compare.pdf}{图3-4　参与规模位次与 OpenRank 位次的对应}{来源：\href{https://github.com/X-lab2017/2026-Global-Open-Source-Development-Report/blob/c51e3441361407b68f10b339d97eff288319b503/data/processed/chapter3/table-3-1-1_global_division_openrank_top30.csv}{[3.1-S1]}同样本重排；统计窗口待核。两轴均为位次，虚线为相同排序；数量字段的单位和账号范围待核，不显示绝对人数，不计算人均效率。}

北京的两项位次为5与14。位次差提示需要进一步核查维护、审阅及跨项目协作角色，现有汇总表尚不能确定具体机制。它比较的是同一截断样本内的排序，不是地区变化或个人能力。

\clearpage
\pdfbookmark[2]{3.1.3 中国省级榜：名次背后的距离}{figure5}
\subsection*{3.1.3 中国省级榜：名次背后的距离}
北京、上海居前，广东、台湾、浙江列第3至第5。广东与台湾相差97.50；浙江约为江苏的2.06倍。相邻名次不代表相近分值。

\plate{figures/31/05_province.pdf}{图3-5　中国省级行政区影响力 Top 30}{来源：固定省级榜 \href{https://github.com/X-lab2017/2026-Global-Open-Source-Development-Report/blob/c51e3441361407b68f10b339d97eff288319b503/data/processed/chapter3/table-3-1-3_china_province_openrank_top30.csv}{[3.1-S3]}；统计窗口待核。条长按本榜榜首归一；全球位次匹配图3-1，横线为未入该Top 30。台湾省、香港特别行政区纳入中国。}

北京、上海也进入全球Top 20。榜次帮助定位，分值呈现距离；选择对标地区时，仍需核对参与基础和项目组合，不能仅按上下相邻名次推断差距成因。

\clearpage
\pdfbookmark[2]{中国城市榜：高位节点与接近位置}{figure6}
\subsection*{中国城市榜：高位节点与接近位置}
北京、上海、杭州、深圳、台北列前五。广州与成都分别为1,921.31和1,918.33，虽分列第6、第7，差值仅2.98。

\plate{figures/31/06_city.pdf}{图3-6　中国城市影响力 Top 30：保留省级归属}{来源：固定城市榜 \href{https://github.com/X-lab2017/2026-Global-Open-Source-Development-Report/blob/c51e3441361407b68f10b339d97eff288319b503/data/processed/chapter3/table-3-1-4_china_city_openrank_top30.csv}{[3.1-S4]}；统计窗口待核。保留原榜次、两位小数及省级归属；条长按本榜榜首归一，两级分值不相加。}

城市榜把地区观察推进到更具体的协作入口。细小分差不足以支持稳健领先的判断；省级和城市分别聚合，不能将城市分值相加还原省级影响。

\clearpage
\pdfbookmark[2]{省市联读：识别多节点组合}{figure7}
\subsection*{省市联读：识别多节点组合}
城市Top 30涉及20个省级地区，其中7个有多个入榜城市。广东、台湾、江苏各有3个，山东、浙江、辽宁、福建各有2个。

\plate{figures/31/07_links.pdf}{图3-7　七个多城市入榜地区的省市关联}{来源：固定省市榜 \href{https://github.com/X-lab2017/2026-Global-Open-Source-Development-Report/blob/c51e3441361407b68f10b339d97eff288319b503/data/processed/chapter3/table-3-1-3_china_province_openrank_top30.csv}{[3.1-S3、S4]}；统计窗口待核。两级分别排序，线条仅表示行政归属，不证明跨城协作。}

广东的深圳、广州进入前六，珠海列第29；江苏的南京、苏州、无锡列第9、11、28。组合信息有助于定位候选合作节点；能否形成跨城协作，还需核查共同项目与实际联系。

\clearpage
\SetRegion{3.2 企业}{violet}{}
\phantomsection\label{ch3-sec32}
\pdfbookmark[1]{3.2 全球开源企业影响力排行榜}{sec32}
\section*{3.2 全球开源企业影响力排行榜}
\noindent\includegraphics[width=\linewidth]{figures/new/nav_32.pdf}\par\vspace{3pt}
微软、华为、谷歌列前三，前两名OpenRank相差25,806.61。Anthropic与OpenAI列第9、第17，显示AI相关企业进入可见的影响位置。

\plate{figures/32/01_global.pdf}{图3-8　全球企业影响力 Top 30：影响、仓库与关联开发者}{来源：固定企业榜 \href{https://github.com/X-lab2017/2026-Global-Open-Source-Development-Report/blob/c51e3441361407b68f10b339d97eff288319b503/data/processed/chapter3/table-3-2-1_global_company_openrank_top30.csv}{[3.2-S1]}；统计窗口、企业归并范围待核。关联开发者非雇员，国别非贡献发生地；各列独立尺度，数量不跨企业相加。}

企业影响需要放回项目组合阅读：关联开发者不等于企业员工，仓库多也不直接表示投入更大。以下分别观察三项指标、榜内集中与跨期位置。

\clearpage
\pdfbookmark[2]{3.2.1 企业画像：三个指标，三种排序}{figure9}
\subsection*{3.2.1 企业画像：三个指标，三种排序}
微软在影响与关联开发者数上居首，华为的仓库数居首。将三个指标在同一组企业中重排，可以辨认参与基础和仓库覆盖的差异。

\plate{figures/32/02_profiles.pdf}{图3-9　全球 Top 30 企业的三指标位次画像}{来源：\href{https://github.com/X-lab2017/2026-Global-Open-Source-Development-Report/blob/c51e3441361407b68f10b339d97eff288319b503/data/processed/chapter3/table-3-2-1_global_company_openrank_top30.csv}{[3.2-S1]}同样本重排；统计窗口待核。条长为（31－位次）/30，只表示30家中的顺序，非原始分值；三列不合成总分。}

Anthropic的影响、仓库与关联开发者位次为9、23、4，IBM为15、5、12。少量仓库可关联较多参与者，多仓库也可能服务不同任务；这些结构不能简化为投入效率。

\clearpage
\pdfbookmark[2]{中国企业榜：头部间距与完整梯次}{figure10}
\subsection*{中国企业榜：头部间距与完整梯次}
华为居首，分值约为阿里巴巴的6.61倍，并超过第二至第五名之和。阿里巴巴与字节跳动分列第二、第三，相差1,818.39。

\plate{figures/32/03_china.pdf}{图3-10　中国企业影响力 Top 30：完整名单与参与结构}{来源：固定中国企业榜 \href{https://github.com/X-lab2017/2026-Global-Open-Source-Development-Report/blob/c51e3441361407b68f10b339d97eff288319b503/data/processed/chapter3/table-3-2-2_china_company_openrank_top30.csv}{[3.2-S2]}；统计窗口和企业归并范围待核。Deepin、openKylin等主体边界沿用来源；标识仅作识别。}

两张Top 30共有四家企业，其三项数值一致。未进入全球Top 30仅表示落在该截断范围之外；比较完整梯次，仍须结合各企业的项目范围。

\clearpage
\pdfbookmark[2]{3.2.2 榜内分布：集中在何处}{figure11}
\subsection*{3.2.2 榜内分布：集中在何处}
全球榜与中国榜前五名分别占各自Top 30分值的51.6\%和79.3\%。同为前列集中，两者的具体结构不同。

\plate{figures/32/04_concentration.pdf}{图3-11　两张 Top 30 榜单的 OpenRank 分值构成}{来源：\href{https://github.com/X-lab2017/2026-Global-Open-Source-Development-Report/blob/c51e3441361407b68f10b339d97eff288319b503/data/processed/chapter3/table-3-2-1_global_company_openrank_top30.csv}{[3.2-S1、S2]}表内复算；统计窗口待核。分母为各榜30家分值之和：925,555.48与261,083.87；两榜有重叠，不相加，不外推整个生态。}

全球榜前两名合计33.3\%，前十名67.2\%；中国榜首位54.1\%，前十名89.7\%。集中位置有助于选择重点调查对象，不能据此给出整个生态的风险等级或市场份额。

\clearpage
\pdfbookmark[2]{3.2.3 月度位置：何时发生换位}{figure12}
\subsection*{3.2.3 月度位置：何时发生换位}
2025年7月至2026年6月，全球月榜首位在3月交换，中国月榜第二名在4月交换。逐月读数保留变化时点。

\plate{figures/32/05_monthly.pdf}{图3-12　全球与中国所选企业的月度位次矩阵}{来源：企业月度表 \href{https://github.com/X-lab2017/2026-Global-Open-Source-Development-Report/blob/c51e3441361407b68f10b339d97eff288319b503/data/processed/chapter3/fig-3-2-2_global_company_monthly_rank.csv}{[3.2-S4、S6]}；2025-07至2026-06，各选10家。格内保留原位次，颜色按相同分档；两榜分别排序，月度名次不平均为年度排名。}

华为在全球月榜3—6月连续第1；字节跳动在中国月榜4—6月连续第2。Anthropic在3—4月达到第6，6月为第13。名次变化只说明相对顺序，不能换算为分值增减或投入效果。

\clearpage
\pdfbookmark[2]{历史位置：进入前列与前列换位}{figure13}
\subsection*{历史位置：进入前列与前列换位}
全球所选企业的历史位次跨越1至309。按企业逐行展开，既保留全部有效读数，也能看见进入前列的阶段。

\plate{figures/32/06_global_history.pdf}{图3-13　全球所选企业历史位次：完整数值矩阵}{来源：全球企业历史表 \href{https://github.com/X-lab2017/2026-Global-Open-Source-Development-Report/blob/c51e3441361407b68f10b339d97eff288319b503/data/processed/chapter3/fig-3-2-1_global_company_yearly_rank.csv}{[3.2-S3]}。保留2016—2026原年度标签，具体起止和2026覆盖期待核；数字为原位次，缺失留灰，不插补。}

微软在所列年度均为第1。华为从2019标签的第66进入2021的第8，2025为第2；英伟达从2024的第12进入2025的第6。历史标签的起止范围仍待核，不能把2026标签当作完整自然年。

\clearpage
\pdfbookmark[2]{中国企业历史：不同阶段的前移}{figure14}
\subsection*{中国企业历史：不同阶段的前移}
阿里巴巴在2016—2020标签下均列第1，华为从2021标签起列第1。字节跳动在2019为第38，2021为第10，2025为第3。

\plate{figures/32/07_china_history.pdf}{图3-14　中国所选企业历史位次：完整数值矩阵}{来源：中国企业历史表 \href{https://github.com/X-lab2017/2026-Global-Open-Source-Development-Report/blob/c51e3441361407b68f10b339d97eff288319b503/data/processed/chapter3/fig-3-2-3_china_company_yearly_rank.csv}{[3.2-S5]}。年份沿用来源标签，各年起止及2026覆盖期待核；只追踪所选10家，缺失留灰，不据此判断全部企业进退。}

PingCAP在所列年度处于第4至第8，飞致云从2016的第21变为2026标签的第8。中间梯次补充了榜首之外的位置变化；名次前移不等于分值或投入等幅增长。

\clearpage
\SetRegion{3.3 基金会}{violet}{}
\phantomsection\label{ch3-sec33}
\pdfbookmark[1]{3.3 全球开源基金会影响力排行榜}{sec33}
\section*{3.3 全球开源基金会影响力排行榜}
\noindent\includegraphics[width=\linewidth]{figures/new/nav_33.pdf}\par\vspace{3pt}
Linux基金会（含子基金会）、开放原子与Apache列前三，榜首分值约为次席的1.96倍。组织位置取决于所关联的项目与仓库范围。

\plate{figures/33/01_board.pdf}{图3-15　全球基金会影响力榜：21个来源条目}{来源：基金会主榜 \href{https://github.com/X-lab2017/2026-Global-Open-Source-Development-Report/blob/c51e3441361407b68f10b339d97eff288319b503/data/processed/chapter3/table-3-3-1_global_foundation_openrank.csv}{[3.3-R1]}。主榜窗口、归并、聚合与人数去重待核；LF与OpenInfra可能重叠，不计算总体份额，不与仓库月均值直接比较。}

本节先保留来源榜，再用LF期末名录与OpenDigger月度数据观察项目组合。两组证据各有范围；项目分值不能直接解释为基金会本身的贡献或治理成效。

\clearpage
\pdfbookmark[2]{3.3.1 月度位置：稳定中的局部换位}{figure16}
\subsection*{3.3.1 月度位置：稳定中的局部换位}
主期内Linux基金会始终第1，NixOS、Rust、Python保持第4、第5、第6。开放原子与Apache则在第2、第3间交替。

\plate{figures/33/02_foundation_monthly.pdf}{图3-16　八家基金会的月度位次变化}{来源：基金会月度表 \href{https://github.com/X-lab2017/2026-Global-Open-Source-Development-Report/blob/c51e3441361407b68f10b339d97eff288319b503/data/processed/chapter3/fig-3-3-1_foundation_monthly_rank.csv}{[3.3-R1]}；2025-07至2026-06。保留8条轨迹；线只连接月度观测，不表示月内变化或分值增减；组织聚合待核。}

开放原子在10个月份列第2，Apache在2025年10月及2026年1月位居其前；两者顺序相邻月份间改变4次。位置稳定与顺序不变并非同一件事，背后的项目组合还需另看。

\clearpage
\pdfbookmark[2]{3.3.2 项目组合：广度与核心集中}{figure17}
\subsection*{3.3.2 项目组合：广度与核心集中}
LF期末名录映射759个GitHub主仓库，其中233个具备主期完整12个月OpenRank。以下比较仅在完整数据样本中展开。

\plate{figures/33/03_portfolios.pdf}{图3-17　四组LF项目组合的可观测结构}{来源：LF期末名录、OpenDigger月度数据 \href{https://github.com/jmertic/lf-landscape/blob/8fcfd93efdf91697b88b97b6ee4abdd2328b3a7b/landscape.yml}{[3.3-R2—R4]}；主期2025-07至2026-06，采集2026-11-05。条带按组内完整样本月均分值；各组覆盖率不同，期末归属不代表全年恒定。}

CNCF组有117个完整仓库，Kubernetes占该组样本分值10.3\%；PyTorch组的5个仓库中，PyTorch与vLLM合计83.2\%。不同结构为维护接续提供调查重点；各组覆盖率不同，不能据此排列整体治理优劣。

\clearpage
\pdfbookmark[2]{重点仓库：把组织名称落到资源}{figure18}
\subsection*{重点仓库：把组织名称落到资源}
233个完整仓库中，PyTorch、Zephyr、vLLM的主期月均OpenRank列前三。完整Top 30同时呈现所属分组和月均值。

\plate{figures/33/04_top30.pdf}{图3-18　LF项目主仓库样本 Top 30}{来源：LF期末名录与OpenDigger \href{https://github.com/jmertic/lf-landscape/blob/8fcfd93efdf91697b88b97b6ee4abdd2328b3a7b/landscape.yml}{[3.3-R2—R4]}；主期2025-07至2026-06。仅含233个完整仓库；月均为12月算术平均，非重算年度网络得分或LF全体排行。}

样本前10个仓库占分值总和40.4\%。这是可观测主仓库的集中程度，不能用来还原基金会主榜总分。

\clearpage
\pdfbookmark[2]{项目月度位置：组织内部仍有变化}{figure19}
\subsection*{项目月度位置：组织内部仍有变化}
在233个完整仓库逐月重排后，PyTorch、Zephyr与vLLM始终列前三；其后项目的相对位置有所改变。

\plate{figures/33/05_repository_monthly.pdf}{图3-19　重点仓库在固定样本中的月度位次}{来源：OpenDigger固定仓库样本 \href{https://open-digger.cn/docs/user_docs/metrics/global_openrank}{[3.3-R3]}；2025-07至2026-06。按233仓逐月重排，同值取最小位次，选8仓展示；期末名录回溯，加入前活动不归因为基金会支持。}

ESLint从第5移至第4，Ceph从第6移至第5；Kubernetes从第4移至第10，Cilium从第8移至第17。稳定的组织位置不能替代对内部项目的观察；后移也不直接等于质量下降或用户流失。

\clearpage
\SetRegion{3.4 开源项目}{violet}{}
\phantomsection\label{ch3-sec34}
\pdfbookmark[1]{3.4 全球开源项目影响力排行榜}{sec34}
\section*{3.4 全球开源项目影响力排行榜}
\noindent\includegraphics[width=\linewidth]{figures/new/nav_34.pdf}\par\vspace{3pt}
OpenHarmony居首，头部对象横跨操作系统、语言工具、云基础设施与AI框架。具体软件与资源集合同时入榜，项目范围需要结合用途阅读。

\plate{figures/34/01_global.pdf}{图3-20　全球项目影响力Top 30的名次与技术用途}{来源：固定全球项目榜 \href{https://github.com/X-lab2017/2026-Global-Open-Source-Development-Report/blob/c51e3441361407b68f10b339d97eff288319b503/data/processed/chapter3/table-3-4-1_global_project_openrank_top30.csv}{[3.4-S1]}。原榜窗口、项目—仓库映射与去重待核；保留30项，不计算人均效率。*Claude Code的公开协作仓库不证明整个产品按开源许可发布。}

从名次继续进入资源、协作与维护，才能理解项目怎样持续供给。本节以OpenHarmony和PyTorch展开观察，并在节末保留两张A3完整画像。

\clearpage
\pdfbookmark[2]{3.4.1 资源结构：名次之外看用途}{figure21}
\subsection*{3.4.1 资源结构：名次之外看用途}
全球Top 30中，AI与机器学习对象占8项（26.7\%）；中国Top 30有13项（43.3\%），另有6项数据与分析资源。

\plate{figures/34/02_atlas.pdf}{图3-21　两份Top 30的项目资源结构}{来源：两份项目榜与同版本用途元数据 \href{https://github.com/X-lab2017/2026-Global-Open-Source-Development-Report/blob/c51e3441361407b68f10b339d97eff288319b503/data/processed/chapter3/table-3-4-1_global_project_openrank_top30.csv}{[3.4-S1、S2]}；原榜窗口待核。单一主用途，每格1项；两榜有4项共同名称，合计56个不同名称。}

语言、编译与环境管理，模型计算与推理，数据库与数据处理，可能处于同一技术系统的不同层次。用途分类帮助识别承担的任务，描述的是入榜对象构成，不是全球技术供给份额。

\clearpage
\pdfbookmark[2]{中国项目榜：多层次的资源供给}{figure22}
\subsection*{中国项目榜：多层次的资源供给}
中国项目榜包括操作系统、计算框架、数据平台与应用工具。OpenHarmony与openEuler居前，CANN、Ascend等提供计算相关资源。

\plate{figures/34/03_china.pdf}{图3-22　中国项目影响力Top 30的名次与技术用途}{来源：固定中国项目榜 \href{https://github.com/X-lab2017/2026-Global-Open-Source-Development-Report/blob/c51e3441361407b68f10b339d97eff288319b503/data/processed/chapter3/table-3-4-3_china_project_openrank_top30.csv}{[3.4-S2]}；原榜窗口、项目映射和去重待核。保留30项，与图3-20同分值尺度；归属沿用来源。}

这是一份独立来源名单，并非从全球Top 30筛选所得。资源能否被其他主体使用，需要继续核对接口、文档、适配和版本支持；名单本身不能概括一国全部技术能力。

\clearpage
\pdfbookmark[2]{位次跃升：选择进一步观察的对象}{figure23}
\subsection*{位次跃升：选择进一步观察的对象}
来源跃升名单保留10项：8项有数值基期，CANN与AgentScope为新进条目。其中7项按主要用途归为AI与机器学习。

\plate{figures/34/08_rising.pdf}{图3-23　中国项目跃升名单的起止位次}{来源：固定跃升名单 \href{https://github.com/X-lab2017/2026-Global-Open-Source-Development-Report/blob/c51e3441361407b68f10b339d97eff288319b503/data/processed/chapter3/table-3-4-4_china_project_rising.csv}{[3.4-S3]}；基期、本期窗口待核。保留全部10项，数值轴1—151；两项新进保留非数值起点，新进不等于首次公开。}

Ascend从第151移至第4，openFuyao从第112移至第21。比较窗口尚未明确，这些记录提供追踪线索，不能称为年度增长或技术突破。解释变化，应先固定项目范围，再核对分值序列、协作与交付。

\clearpage
\pdfbookmark[2]{3.4.2 资源与交付：从清单到版本}{resources-delivery}
\subsection*{3.4.2 资源与交付：从清单到版本}
OpenHarmony期末默认源码清单含512项，其中第三方依赖183项、系统框架117项、基础能力109项，提供分解维护责任的起点。

\plate{figures/34/portrait_oh_resources.pdf}{图3-24　OpenHarmony默认源码清单的资源构成}{来源：2026-06-30固定manifest \href{https://github.com/openharmony/manifest/tree/36a80ca0ea0101ed0df1d7b773559e428f6eef0f}{[3.4-S20]}，按本地路径首段计数。面积为条目数，非代码量、原创性或自主性；默认清单不是全社区仓库总表或榜单映射。}

PyTorch核心仓库主期新增23,777条PR、7,329条Issue，发布7个非预发布版本；后半期PR较前半期增加31.4\%，应进一步核对评审与兼容维护容量。

\plate{figures/34/portrait_delivery.pdf}{图3-25　PyTorch核心仓库的工作流入与发布节奏}{来源：GitHub创建及发布记录 \href{https://api.github.com/search/issues?q=repo:pytorch/pytorch+is:pr+created:2025-07-01..2026-06-30&per_page=1}{[3.4-S11、S12]}；2025-07-01至2026-06-30（UTC），月和与整期一致。发布按published\_at及prerelease=false；记录含自动化与堆叠工作流，不等同独立工作量。}

\clearpage
\pdfbookmark[2]{3.4.3 贡献接续：平台转换与持续参与}{people-continuity}
\subsection*{3.4.3 贡献接续：平台转换与持续参与}
OpenHarmony主期贡献榜并集有513个平台账号标识；11月起前100条目均来自AtomGit。平台构成转换不等于人员迁移完成。

\plate{figures/34/portrait_oh_platform.pdf}{图3-26　OpenHarmony贡献榜样本的平台构成}{来源：OpenDigger社区贡献榜 \href{https://oss.open-digger.cn/companies/meta/pytorch/community_openrank_details.json}{[3.4-S15]}；2025-07至2026-06，每月前100条。按平台+稳定账号ID识别，跨平台未归并；只反映榜内构成。}

PyTorch项目标签榜形成244个账号标识，61个出现9—12个月（25.0\%），105个只出现1—2个月（43.0\%）。持续与阶段性参与并存。

\plate{figures/34/portrait_pytorch_people.pdf}{图3-27　PyTorch贡献榜名次、当期贡献值与持续性}{来源：OpenDigger项目标签贡献榜、固定维护者名册 \href{https://oss.open-digger.cn/companies/meta/pytorch/community_openrank_details.json}{[3.4-S15、S18]}；主期12个月。热图追踪6月前12条，灰格为未入前100；已知机器人规则无命中，不代表全部人工确认。}

将贡献榜与固定维护者名册联读，才能进一步检查模块责任、共同维护及替补安排；账号进入榜单本身不证明拥有维护权限。

\clearpage
\pdfbookmark[2]{跨项目共同出现：定位协作线索}{figure28}
\subsection*{跨项目共同出现：定位协作线索}
六个代表仓库的月度高活跃账号样本中，vLLM与SGLang有36个共同账号，PyTorch与vLLM有22个。公开参与记录为跨项目关系提供了调查入口。

\plate{figures/34/06_connections.pdf}{图3-28　六个代表仓库的高活跃账号交集}{来源：六仓activity\_details \href{https://oss.open-digger.cn/github/pytorch/pytorch/activity_details.json}{[3.4-S6]}；2025-07至2026-06。每仓每月最多100条，先截榜后剔除已知机器人，不回补；归一登录名匹配，改名与稳定身份未完整处理。}

vLLM与SGLang的36个共同账号中，21个曾在同月进入两份名单；LangChain与Dify的16个共同账号中，15个有同月记录。时间重合缩小了调查范围，但仍需具体Issue、PR或适配记录确认任务关系。

\clearpage
\pdfbookmark[2]{出现频次：持续入样与维护职责}{figure29}
\subsection*{出现频次：持续入样与维护职责}
PyTorch核心仓库有36个账号进入9—12个月高活跃名单，Kubernetes为30个，vLLM为25个。它们提供持续参与的线索，尚不能直接认定为拥有维护权限的维护者。

\plate{figures/34/07_frequency.pdf}{图3-29　高活跃账号样本的月度出现频次}{来源：六仓月度高活跃样本 \href{https://oss.open-digger.cn/github/pytorch/pytorch/activity_details.json}{[3.4-S6]}；2025-07至2026-06。分母为各仓12月账号并集，非连续参与或留存率。此处PyTorch核心仓库36个，与图3-27项目标签OpenRank榜61个范围不同。}

六个仓库均以进入1—2个月名单的账号最多，可能涉及阶段性任务、问题反馈或截榜机制。评估接续能力还需追踪评审、发布及问题处理责任，以及关键任务是否有替补；名单频次不能替代全社区留存率。

\clearpage
\pdfbookmark[2]{3.4.4 复用与维护：依赖约束与责任接续}{reuse-maintenance}
\subsection*{3.4.4 复用与维护：依赖约束与责任接续}
期末vLLM与SGLang均声明PyTorch及Transformers依赖，具体版本与启用方式不同。使用方需要核对兼容条件和支持期限。

\plate{figures/34/05_dependencies.pdf}{图3-30　两个推理项目的直接软件依赖}{来源：不晚于2026-06-30的固定依赖文件 \href{https://github.com/vllm-project/vllm/tree/ab80b3dff4829768347b322c7a5d444e522ef7b3}{[3.4-S5、S22]}（同组证据）及PyPI版本制品 \href{https://pypi.org/pypi/torch/2.12.1/json}{[3.4-S16]}。箭头为使用方→依赖项；仅三组选定依赖，非完整供应链或实际安装状态。}

\plate{figures/34/portrait_maintenance.pdf}{图3-31　PyTorch维护观测的分布与缺失月份}{来源：OpenDigger维护数据及固定治理文件 \href{https://github.com/X-lab2017/open-digger/blob/master/src/metrics/chaoss.ts}{[3.4-S10、S18]}。响应与Issue年龄8/12月，PR解决4/12月；未响应代入15天，关闭也算响应。缺月留白，月度分位数不平均为年度值。}



\clearpage\phantomsection\label{ch3-portraits}
\pdfbookmark[2]{完整项目画像：PyTorch与OpenHarmony}{portraits}
\includepdf[pages=-,fitpaper=true,pagecommand={\thispagestyle{empty}}]{atlas/atlas.pdf}

\clearpage
\SetRegion{3.5 新势力}{violet}{}
\phantomsection\label{ch3-sec35}
\pdfbookmark[1]{3.5 全球新势力开源仓库与项目排行榜}{sec35}
\section*{3.5 全球新势力开源仓库与项目排行榜}
\noindent\includegraphics[width=\linewidth]{figures/new/nav_35.pdf}\par\vspace{3pt}
项目榜呈现整体供给，仓库榜揭示具体组件。CANN居项目榜首，Hermes Agent、DeepAgents等工具与框架随后；30项中有20个单仓项目、10个多仓项目。

\plate{figures/35/01_projects.pdf}{图3-34　新项目 Top 30 的协作影响与仓库范围}{来源：固定新项目榜及身份映射 \href{https://github.com/X-lab2017/2026-Global-Open-Source-Development-Report/tree/c51e3441361407b68f10b339d97eff288319b503/data/processed/chapter3}{[3.5-E1、E2、E5]}；沿用来源主期2025-07至2026-06，完整纳入范围和新建资格待核。官方标识采集于2026年11月。}

两份榜单的平台覆盖与纳入规则不同。联读它们，可以同时观察智能体工具的项目集聚与基础软件的组件分工。

\clearpage
\pdfbookmark[2]{3.5.1 榜单全景：进入具体协作入口}{figure35}
\subsection*{3.5.1 榜单全景：进入具体协作入口}
仓库榜以OpenClaw居首，并列出算子库、框架适配和系统组件。完整路径保留平台与命名空间，便于识别同一软件栈内部的分工。

\plate{figures/35/02_repos.pdf}{图3-35　新仓库 Top 30 的完整路径与协作影响}{来源：固定新仓库榜、OpenDigger回算 \href{https://github.com/X-lab2017/2026-Global-Open-Source-Development-Report/tree/c51e3441361407b68f10b339d97eff288319b503/data/processed/chapter3}{[3.5-E1—E3、E5]}。G为GitHub，A为AtomGit；30仓主期回算最大绝对差0.02。人数不跨仓相加；平台新建不等于技术首次公开。}

cann与Ascend各有8席，不能视为16个独立新项目，也不能全部并入CANN。

\clearpage
\pdfbookmark[2]{3.5.2 技术构成：单位与平台共同改变}{figure36}
\subsection*{3.5.2 技术构成：单位与平台共同改变}
新仓库名单中，AI基础软件占18项（60.0\%）；新项目名单中，智能体与Agent工具占25项（83.3\%）。两者呈现不同供给结构。

\plate{figures/35/03_macro.pdf}{图3-36　两份 Top 30 的技术构成与平台范围}{来源：固定榜表与编辑分类 \href{https://github.com/X-lab2017/2026-Global-Open-Source-Development-Report/tree/c51e3441361407b68f10b339d97eff288319b503/data/processed/chapter3}{[3.5-E1、E2]}。每项一个主领域，等面积为数量而非分值；语言仅仓库榜汇总，TS、C\#、QML各1项；分类版本记录随源码保留。}

仓库榜含20项AtomGit与10项GitHub，项目榜为1项AtomGit与29项GitHub。统计单位和平台范围同时改变，不能把技术差异全部归因于仓库拆分。

仓库语言以Python与C++为主，分别15项、12项，合计90\%。该汇总不提供逐仓语言与用途的交叉对应，也不代表全球技术份额。

\clearpage
\pdfbookmark[2]{3.5.3 CANN：多仓软件栈的组件分工}{figure37}
\subsection*{3.5.3 CANN：多仓软件栈的组件分工}
来源项目记录纳入CANN的68个仓库，其中8个经官方关系核对的组件进入新仓库Top 30。局部窗口呈现算子、编程工具和运行支撑的分工。

\plate{figures/35/04_cann.pdf}{图3-37　CANN项目范围与八个入榜组件的功能构成}{来源：固定榜表与官方组件关系 \href{https://github.com/X-lab2017/2026-Global-Open-Source-Development-Report/tree/c51e3441361407b68f10b339d97eff288319b503/data/processed/chapter3}{[3.5-E1、E2]}；开放进程见华为公告 \href{https://www.huawei.com/cn/news/2025/8/ascend-summit-cann-open-source}{[3.5-E4]}。顶部为68仓项目指标，面积仅8个入榜组件，分母10,028.22；非68仓全景或项目总分分解。}

8个组件分值中，ops-transformer、ops-nn、ops-math合计58.3\%，编程工具21.0\%，运行支撑20.7\%。组件供给能否持续复用，需要版本配套、接口兼容与维护责任共同支持。

2025年8月5日全面开源开放公告与9月20日技术指导委员会公告记录了建设进程，不单独证明全部计划已交付，也不解释榜单分值的成因。

\clearpage
\SetRegion{3.6 开源模型}{violet}{}
\phantomsection\label{ch3-sec36}
\pdfbookmark[1]{3.6 全球开源模型影响力排行榜}{sec36}
\section*{3.6 全球开源模型影响力排行榜}
\noindent\includegraphics[width=\linewidth]{figures/new/nav_36.pdf}\par\vspace{3pt}
本节采用2026年6月ModelRank月榜，观察主期末的模型版本、任务与发布账户。前十名含图像生成4席、文本生成5席、机器人1席。

\plate{figures/36/01_models.pdf}{图3-38　模型影响力前十名与复用相关记录}{来源：ModelRank及模型元数据 \href{https://github.com/X-lab2017/2026-Global-Open-Source-Development-Report/blob/c51e3441361407b68f10b339d97eff288319b503/preview/content/gosr/model-monthly-source/modelrank-2026-06-top100.raw.json}{[3.6-E1—E3]}；2026年6月月榜。下载、衍生的窗口和去重待核，不作为当月新增；标识为当前官方头像。ModelRank不与OpenRank直接比较。}

Qwen2.5-7B基础版的影响位次高于指令版，下载与衍生记录则相反。系列名称相近不能替代具体版本识别；文件访问量也不能替代复用网络中的位置。

FLUX.1-dev当前模型卡标注非商业许可，提示复用影响与开放条件须分开判断。该榜是版本观察样本，不是年度累计或完整开源认证清单。

\clearpage
\pdfbookmark[2]{3.6.1 任务构成：供给广度与头部位置}{figure39}
\subsection*{3.6.1 任务构成：供给广度与头部位置}
文本生成与图文理解分别有38个和20个版本，合计58\%。图像生成只有6个版本，却有5个进入前20；唯一入榜的机器人模型smolvla\_base列第9。

\plate{figures/36/02_tasks.pdf}{图3-39　任务入榜数量与名次分布}{来源：ModelRank原任务字段 \href{https://github.com/X-lab2017/2026-Global-Open-Source-Development-Report/blob/c51e3441361407b68f10b339d97eff288319b503/preview/content/gosr/model-monthly-source/modelrank-2026-06-top100.raw.json}{[3.6-E1]}；2026年6月。橙色为前20子集，短线保留名次；5条无任务标签，保留缺失，不用当前模型卡回填。}

图像生成较集中于前列，文本生成横跨第3至第98。语音、检索和表示学习等任务也提供资源入口，席位较少不代表实际需求较低。综合榜用于发现候选，应用选择还需任务评测、适配成本和维护证据。

\clearpage
\pdfbookmark[2]{3.6.2 发布账户：规模与布局分别观察}{figure40}
\subsection*{3.6.2 发布账户：规模与布局分别观察}
100个版本来自22个账户。Qwen与google合计56席，而前20名分布于9个账户。数量、头部席位和任务布局反映不同侧面。

\plate{figures/36/03_publishers.pdf}{图3-40　发布账户的入榜数量与任务布局}{来源：ModelRank与账户身份记录 \href{https://github.com/X-lab2017/2026-Global-Open-Source-Development-Report/blob/c51e3441361407b68f10b339d97eff288319b503/preview/content/gosr/model-monthly-source/modelrank-2026-06-top100.raw.json}{[3.6-E1、E3]}；2026年6月。22账户分别统计，不并入母公司；21个官方头像，缺头像者用文字；任务数0表示标签未取得。}

Qwen的36席中9席进入前20，google的20席中2席进入前20；krea的两席均在前20。账户不等于模型家族或母公司，版本数量与头部位置也不能合成为组织能力总分。